\documentclass[12pt,a4paper]{article}
\usepackage[utf8]{inputenc}
\usepackage[margin=2.5cm]{geometry}
\usepackage[numbers]{natbib}
\bibliographystyle{ieeetr}

\title{PE5 -- Vista previa individual\\Sección 1: Introducción}
\author{Nieves Sánchez, Jimmy Samuel}
\date{}

\begin{document}
\maketitle



\section{Introducción}

\textbf{Problema.} Los procesos de adopción y cruza responsable de mascotas se
desarrollan hoy, en su mayoría, de manera informal: mediante publicaciones dispersas
en redes sociales, grupos de mensajería o acuerdos verbales entre particulares, sin
un canal centralizado que garantice la veracidad de la información publicada ni la
trazabilidad del proceso. La elicitación realizada en PE1--PE2 documentó esta
problemática de forma directa: los propietarios y adoptantes entrevistados señalaron
reiteradamente la falta de transparencia de algunos refugios sobre el estado
sanitario de las mascotas muchas veces por saturación y falta de financiamiento
fijo, casos de mascotas que no correspondían a lo anunciado, y estafas en
adopciones en línea mediante cobros sucesivos por concepto de envío o implementos. A esto se suma la ausencia de mecanismos de verificación de la
identidad del propietario y de la calidad de las fotografías publicadas, lo que
facilita la proliferación de perfiles fraudulentos o irrelevantes. MundiPets
responde a este problema como una plataforma digital que centraliza la información
de las mascotas, valida su publicación y ofrece mecanismos de apoyo incluyendo
componentes basados en inteligencia artificial para moderar contenido y sugerir
compatibilidad, sin sustituir en ningún caso la evaluación profesional veterinaria.

\textbf{Contexto.} El dominio de MundiPets involucra a propietarios de mascotas,
adoptantes, interesados en procesos de cruza, refugios y veterinarias aliadas, cada
uno con necesidades y niveles de acceso distintos dentro de la plataforma. El
entorno organizacional es el de un proyecto académico integrador desarrollado
durante las prácticas PE1 a PE5 de la asignatura Ingeniería de Requisitos
(ISR-401), en el que el equipo actuó simultáneamente como analista de requisitos y
como validador de su propio trabajo mediante inspecciones formales. La evidencia
recogida en PE1--PE2 (entrevistas y cuestionarios a propietarios, adoptantes,
interesados en cruza y refugios) reveló, además, posturas heterogéneas frente a la
privacidad de la información: mientras una parte de los usuarios entrevistados
prefiere compartir abiertamente el historial médico de sus mascotas porque considera
que beneficia al animal ante una emergencia, otra parte delimita como sensible
únicamente el tipo de sangre y las enfermedades, reservando ese dato al veterinario. Esta heterogeneidad condicionó directamente las decisiones de
especificación registradas en el ERS y sustenta, en particular, los requisitos de
base legal y clasificación de riesgo de los componentes de IA descritos en la
Sección 7.

\textbf{Objetivos.} El objetivo general de la PE5 es consolidar e integrar el
trabajo de las cuatro prácticas previas en un informe final único que demuestre,
con evidencia auditable, que el ERS de MundiPets cumple los atributos de calidad
exigidos por \citep{ieee29148}. De este objetivo general se desprenden cinco
objetivos específicos, correspondientes a las cinco secciones centrales del
informe: (1)~auditar la calidad del ERS mediante las seis métricas M1--M6 con
aritmética visible y trazable a los conteos base; (2)~cerrar la trazabilidad
end-to-end del sistema, resolviendo requisitos huérfanos, artefactos huérfanos y
cadenas rotas; (3)~especificar formalmente los componentes de inteligencia
artificial del sistema, con requisitos funcionales y no funcionales medibles;
(4)~consolidar una versión final única y coherente del ERS/SRS, con control de
versión e identificadores verificados; y (5)~preparar y sustentar la defensa
técnica individual ante el tribunal, evidenciando dominio del proyecto completo.

\textbf{Estructura del documento.} Tras este Fundamento Teórico, que reúne las
bases conceptuales y normativas aplicadas por cada integrante siguiendo los
principios de elicitación y validación de \citep{pohl2025}, el informe se organiza
en diez secciones: la Sección~2 describe la metodología de ingeniería de
requisitos efectivamente seguida en PE1--PE5; la Sección~3 presenta el ERS/SRS
final conforme a la estructura de \citep{ieee29148}; la Sección~4 ofrece una
lectura interpretativa de los modelos UML vigentes; la Sección~5 documenta la
inspección formal y la reinspección de defectos; la Sección~6 desarrolla la
gestión y la trazabilidad, incluida la matriz final; la Sección~7 especifica los
requisitos de los componentes de inteligencia artificial; la Sección~8 presenta
el instrumento y los resultados de la auditoría de calidad del ERS; la
Sección~9 recoge la retrospectiva del equipo; y la Sección~10 cierra con cinco
conclusiones integradoras, una por cada unidad del proyecto. El Apéndice reúne el
instrumento de auditoría completo, el banco de preguntas del tribunal, la
declaración individual de aporte y la matriz de trazabilidad final en su versión
legible.

\bibliography{intro_referencias}

\end{document}
